\section{Ограничения}
\label{sec:limits}

\paragraph{Один прогон RL, один seed.} Все числа, связанные с обучением с
подкреплением, получены в единственном прогоне с \code{seed=42}: 64 шага
оптимизатора, 1024 промпта, группа из 8 роллаутов. Этого достаточно, чтобы
зафиксировать \emph{направление} эффекта и его величину на порядок, но не
для оценки вариативности обучения. Парный бутстрэп по пользователям
покрывает только разброс между пользователями при фиксированной обученной
политике: он не говорит, повторится ли падение при другом seed, другом
порядке промптов или другом чекпоинте базы. Поэтому мы формулируем вывод
как «в этой постановке RL снижает релевантность и повышает награду», а не
как универсальную цену RL; три seed и перебор порядка промптов ---
обязательный минимум перед любым более сильным утверждением.

\paragraph{Две разные базовые модели.} RL обучен поверх чекпоинта единого
прогона (3\,000 шагов, тестовый NDCG@10 0.0282); отдельный прогон базы на
30\,000 шагов даёт 0.0623, но RL на нём не запускался. Поэтому: (а) все
проценты падения в работе считаются относительно 0.0282; (б) неизвестно,
насколько абсолютный размер падения зависит от силы опорной политики;
(в) цифры отдельного прогона базы не следует сравнивать с RL как с
«до и после». Подъём базы с 3\,000 до 30\,000 шагов дал $+120\%$ по
NDCG@10 --- при таком масштабе разницы бюджета обучения эффект RL,
измеренный на коротком чекпоинте, нельзя переносить на длинную.

\paragraph{Reward model обучена на логах периода обучения.}
Распределение её входов на сгенерированных айтемах отличается от
обучающего: это классическая причина переоптимизации, и здесь она не
устранена --- обучение на данных, порождённых политикой, в работе не
проводилось.

\paragraph{Разомкнутый симулятор.} Пять сгенерированных айтемов не влияют
на историю и на последующие награды. Поэтому выводы о долгосрочных целях
(H5) в этой постановке невозможны: «возврат в течение 7 дней» ---
статическая оценка, а не результат последовательности действий.

\paragraph{Границы измерений.}
\begin{itemize}[leftmargin=1.2em,itemsep=1pt,topsep=2pt]
  \item \code{played\_ratio\_pct} достигает 159\% при повторном
        прослушивании, а сигмоида ограничена единицей: цель регрессии
        клипуется, и модель систематически недооценивает очень высокие
        значения;
  \item MAP@$k$ нормируется на $\min(|\text{цели}|,k)$, поэтому MAP@100
        ниже MAP@10; сравнение с литературой по MAP ограничено
        одним значением $k$;
  \item в данных есть дубликаты прослушиваний (1.24\% событий обучающего
        периода совпадают по паре uid/айтем/время), дедупликация не
        применяется --- как и в официальном протоколе;
  \item официальный \code{timesplit} включает события ровно на
        $t=26\,000\,000$ в тест, а наш полуинтервал их исключает; это
        три события из 157\,696 целей и не влияет на вывод.
\end{itemize}

\paragraph{Сравнение с SASRec.} Опубликованные значения SASRec приведены
как внешняя точка отсчёта. У них другой протокол препроцессинга, другой
бюджет обучения и неизвестные нам детали инициализации и подбора
гиперпараметров, поэтому отношение «в 1.20 раза» характеризует масштаб, а не
доказанное отставание. Ни один вывод работы не опирается на это сравнение.

\paragraph{То, что выходит за рамки работы.} Абляции награды и KL
($\beta_{\text{KL}} \in \{0, 0.01, 0.2\}$, награды like/blend, DPO,
filtered BC), режим «только органические прослушивания», воспроизводимость
на других seed и перебор гиперпараметров. Соответствующие стадии
реализованы (\code{ablate}) и запускаются одной командой.
